\chapter{群}
\orig{69--88}

\section{群的定义}

\subsection*{1. 群的两个定义}

\textbf{群的第一定义。} 一个非空集合 $G$ 对于一个叫做乘法的代数运算来说作成一个群，如果：
\begin{enumerate}[label=\Roman*.,leftmargin=3em]
\item $G$ 对这个乘法是闭的；
\item 乘法适合结合律；
\item 对 $G$ 的任意 $a,b$，方程
\[
ax=b,\qquad ya=b
\]
在 $G$ 中都有解。
\end{enumerate}

\textbf{群的第二定义。} 一个非空集合 $G$ 对乘法作成一个群，如果满足闭性、结合律，并且：
\begin{enumerate}[label=\Roman*.,leftmargin=3em,start=4]
\item $G$ 中存在一个固定的左单位元 $e$，使 $ea=a$（$\forall a\in G$）；
\item 对每个 $a\in G$，存在左逆元 $a^{-1}$，使
\[
a^{-1}a=e.
\]
\end{enumerate}

这里第二定义中“固定的左单位元”这一点不能随意删去。比如令 $S=\{a,b\}$，乘法表为
\[
\begin{array}{c|cc}
 &a&b\\\hline
 a&a&b\\
 b&a&b
\end{array}
\]
则方程 $yb=a$ 在 $S$ 中无解，所以 $S$ 不是群。

如果把第二定义的最后条件改成：对每个 $a$，允许另找一个依赖于 $a$ 的左单位元 $e'$，并有某个 $a^{-1}$ 满足 $a^{-1}a=e'$，仍不足以保证成群。上面的 $S$ 就满足这种弱化条件：$a,b$ 都可分别充当相应元素的左单位元，但仍不是群。

\subsection*{2. 有限群的定义}

对于有限非空集合 $G$，还可以用闭性、结合律和消去律来刻画群：若
\[
ax=ax'\Longrightarrow x=x',
\qquad
ya=y'a\Longrightarrow y=y',
\]
则 $G$ 为群。这里“有限”不能省略。例如
\[
G=\mathbb Z\setminus\{0\}
\]
在普通乘法下满足闭性、结合律与消去律，却不是群，因为一般整数没有乘法逆元。

\subsection*{3. 群元的阶}

\textbf{定义。} 设 $a\in G$。若存在正整数 $m$ 使
\[
a^m=e,
\]
则满足它的最小正整数 $m$ 叫做 $a$ 的\textbf{阶}；若不存在这样的 $m$，则说 $a$ 为无限阶元素。

有限群的每个元素都有有限阶，但“每个元素都有有限阶”并不能推出群有限。

例如由 $1$ 的所有复数单位根组成的集合，在复数乘法下构成群。若
\[
\varepsilon^n=1,
\qquad
\eta^m=1,
\]
则
\[
(\varepsilon\eta)^{mn}=1.
\]
当 $(m,n)=1$ 时，$\varepsilon\eta$ 是一个 $mn$ 次单位根；一般若 $(m,n)=d$，则它至少是一个 $mn/d$ 次单位根。并且
\[
(\varepsilon^{-1})^n=(\varepsilon^n)^{-1}=1.
\]
因此全体单位根构成一个无限群，而其中每个元素阶都有限。

另一方面，元素全为无限阶的群不可能存在，因为任意群的单位元 $e$ 的阶总是 $1$。除了单位元以外全为无限阶元素的群则很多，例如整数加群；在有理数乘群 $\mathbb Q^*$ 中，$1$ 的阶为 $1$，$-1$ 的阶为 $2$，其余元素均为无限阶。

再给两个“无限群中所有元素阶有限”的例子：

\textbf{例 1。} 设 $\mathbb Z_n$ 为模 $n$ 的剩余类环，令 $G=\mathbb Z_n[x]$ 的全体多项式在加法下成群。$G$ 有无限多个元素，但每个元素的加法阶都有限（并且都整除 $n$）。

\textbf{例 2。} 特征为素数 $p$ 的无限域，在加法下构成无限群，而每个非零元素的阶都是 $p$。

\section{群的同态}

\textbf{1.} 若群 $G$ 满射同态到一个代数系统 $\bar G$，则 $\bar G$ 自动是群。反过来，一个代数系统满射同态到群，却未必本身成群。

\textbf{例 1。} 令 $\bar G$ 为全体奇整数，在普通乘法下并不成群；令 $G=\{e\}$ 为一元素群。常值映射
\[
\phi:\bar a\longmapsto e
\]
是 $\bar G$ 到 $G$ 的满射同态，但 $\bar G$ 不是群。

\textbf{例 2。} 令 $\bar G=\mathbb N$ 为全体自然数，在普通加法下不成群；令 $G=\{1,-1\}$ 在普通乘法下成群。定义
\[
\phi(a)=\begin{cases}
1,&a\text{ 为偶数},\\
-1,&a\text{ 为奇数}.
\end{cases}
\]
则 $\bar G$ 满射同态到 $G$，但 $\bar G$ 仍不是群。

\textbf{2.} 若群 $G$ 满射同态到群 $\bar G$，则 $G$ 的单位元 $e$ 的像是 $\bar G$ 的单位元，且 $a^{-1}$ 的像是 $a$ 的像的逆元。反过来，单位元的逆像未必只有 $G$ 的单位元；两个互逆元素的某些逆像也未必互逆。

例如令 $G=\mathbb Z$ 为整数加群，令
\[
\bar G=\{e,b,c\},
\qquad
\begin{array}{c|ccc}
 &e&b&c\\\hline
 e&e&b&c\\
 b&b&c&e\\
 c&c&e&b
\end{array}
\]
并定义
\[
\phi(x)=\begin{cases}
e,&x\equiv0\pmod3,\\
b,&x\equiv1\pmod3,\\
c,&x\equiv2\pmod3.
\end{cases}
\]
则 $G$ 与 $\bar G$ 同态。可是 $3$ 也是 $e$ 的逆像，而 $3$ 不是 $G$ 的单位元；又 $b,c$ 互逆，但例如它们的逆像 $4,5$ 并不互逆（在加法群中 $4+5\ne0$）。

\textbf{3.} 在群的同态映射下，对应元素的阶不一定相同。

仍取上例，$3,4,5\in\mathbb Z$ 都是无限阶元素，而
\[
3\mapsto e,\qquad4\mapsto b,\qquad5\mapsto c,
\]
其中 $e$ 的阶为 $1$，$b,c$ 的阶为 $3$。

又例如令
\[
G=\left\{1,\frac{-1+\sqrt3i}{2},\frac{-1-\sqrt3i}{2}\right\},
\qquad
\bar G=\{1\},
\]
都取普通乘法，并令 $\phi(x)=1$。则 $G$ 满射同态到 $\bar G$，但两个非平凡三次单位根的阶为 $3$，而 $1$ 的阶为 $1$。

\textbf{4.} 群的自同态满射未必是自同构。例如 $\mathbb R[x]$ 在加法下是群，
\[
\phi:f(x)\longmapsto f'(x)
\]
是一个满射自同态，但不是一一映射，因为不同常数多项式都有导数 $0$。

\section{几个具体群}

\subsection*{1. 变换群}

\textbf{定义。} 一个集合 $A$ 的若干个一一变换对变换复合作成的群，叫做 $A$ 的一个\textbf{变换群}。

\textbf{定理。} 设 $G$ 是集合 $A$ 的若干变换组成的集合，且在变换复合下本身成群。如果 $G$ 中含有 $A$ 的恒等变换 $e$，那么 $G$ 一定是变换群。

“含恒等变换”这一点不能直接忽略。例如 $A=\{1,2\}$，令
\[
\tau:1\mapsto1,\qquad2\mapsto1,
\]
并取 $G=\{\tau\}$。有 $\tau^2=\tau$，因此抽象地说 $G$ 是一元素群，但 $\tau$ 不是一一变换，故 $G$ 不是变换群。

不过，“含恒等变换”可由较弱条件替代：若 $G$ 在复合下成群，并且 $G$ 至少含一个满射变换或至少含一个单射变换，那么 $G$ 就是变换群。

证明满射情形。设 $\tau\in G$ 为满射，$\varepsilon$ 是抽象群 $G$ 的单位元。由
\[
(a^\tau)^\varepsilon=a^{\tau\varepsilon}=a^\tau
\]
并利用 $\tau$ 的满射性，得到 $\varepsilon$ 在 $A$ 上就是恒等变换。于是每个群元的抽象逆元都是其真正的逆变换，因此所有元素均是一一变换。

单射情形类似：设 $\lambda\in G$ 单射，由
\[
(a^\varepsilon)^\lambda=a^{\varepsilon\lambda}=a^\lambda
\]
和单射性得到 $a^\varepsilon=a$，故 $\varepsilon$ 仍是恒等变换。

但一个非一一变换仍可能有\emph{单侧}逆。取
\[
A=\{1,2,3,\ldots\},
\]
定义
\[
\tau(1)=1,\quad\tau(2)=1,\quad\tau(3)=2,\quad\tau(4)=3,\ldots
\]
以及
\[
\sigma(1)=1,\quad\sigma(2)=3,\quad\sigma(3)=4,\quad\sigma(4)=5,\ldots.
\]
则 $\tau$ 不是一一变换，但按本节采用的复合记号有 $\sigma\tau=e$。这说明单侧逆不足以推出双射。

\subsection*{2. 置换群}

\textbf{定义。} 有限集合的一个一一变换叫做一个\textbf{置换}。有限集合的若干置换组成的群叫做一个\textbf{置换群}。含 $n$ 个元素的集合的全体置换组成的群叫做 $n$ 次\textbf{对称群}，记为 $S_n$。

一个 $n$ 元集合的置换称为 $k$--循环置换，如果它把
\[
a_{i_1}\mapsto a_{i_2}\mapsto\cdots\mapsto a_{i_k}\mapsto a_{i_1}
\]
而使其他元素不变。

一个置换不一定本身是一个循环。例如
\[
\begin{pmatrix}
1&2&3&4\\
2&1&4&3
\end{pmatrix}
=(12)(34)
\]
不是单一循环。

每一个 $n$ 元置换都能写成若干互不相交循环的乘积；若允许相交循环，表示一般不唯一。例如
\[
(25)=(12)(15)(12)=(15)(12)(15)
=(45)(34)(23)(34)(45).
\]

$S_3$ 有六个元素
\[
S_3=\{(1),(12),(13),(23),(123),(132)\},
\]
是非交换群，并且是阶数最小的非交换群。

\subsection*{3. 循环群}

\textbf{定义。} 若群 $G$ 的每个元素都是某个固定元 $a$ 的幂，则称 $G$ 为\textbf{循环群}，并说 $G$ 由 $a$ 生成，记作
\[
G=(a).
\]
$a$ 叫做一个生成元。

\textbf{1)} 写成 $G=(a^m)$ 时，$a^m$ 不一定是 $(a)$ 中 $a$ 的最小正方幂。比如六元循环群中
\[
(a)=(a^5).
\]

\textbf{2)} 循环群的同态满射像仍是循环群；反之不真。例如符号映射
\[
\phi:S_3\longrightarrow\{1,-1\},
\qquad
\phi(\pi)=\begin{cases}1,&\pi\text{ 为偶置换},\\-1,&\pi\text{ 为奇置换},\end{cases}
\]
是满射同态，目标群是循环群，而 $S_3$ 非交换，故不是循环群。

\textbf{3)} 循环群的真子群一定是循环群，反之不真。例如 Klein 四元群
\[
G=\{e,a,b,c\},
\qquad
\begin{array}{c|cccc}
 &e&a&b&c\\\hline
 e&e&a&b&c\\
 a&a&e&c&b\\
 b&b&c&e&a\\
 c&c&b&a&e
\end{array}
\]
不是循环群，但它的三个非平凡真子群
\[
\{e,a\},\quad\{e,b\},\quad\{e,c\}
\]
都是二元循环群。可具体实现为
\[
e=(1),\quad a=(12)(34),\quad b=(13)(24),\quad c=(14)(23).
\]

\textbf{4)} 阶为 $p^n$（$p$ 为素数）的循环群恰有 $n-1$ 个非平凡真子群；若去掉“循环群”条件则不真。上面的 Klein 四元群阶为 $2^2$，却有三个非平凡真子群，而不是 $2-1=1$ 个。

\textbf{5)} 一个循环群也可以由多于一个元素共同生成。例如整数加群满足
\[
\mathbb Z=(1),
\qquad
\mathbb Z=(2,3).
\]

\section{子群}

\subsection*{1. 子群}

\textbf{定义。} 群 $G$ 的一个子集 $H$ 若对 $G$ 的乘法本身也作成群，则称 $H$ 为 $G$ 的一个\textbf{子群}。

这里必须强调：$H$ 要沿用 $G$ 的运算，而不能只要求 $H$ 在某个另外的运算下成群。例如 $\mathbb Q$ 是有理数加群，而正有理数集合 $\mathbb Q_{>0}$ 在普通乘法下虽然成群，却不是 $\mathbb Q$ 作为加法群的子群。

构成子群的闭性与取逆两个条件彼此独立。非空子集 $H\subseteq G$ 是子群，当且仅当
\[
a,b\in H\Longrightarrow ab\in H,
\qquad
 a\in H\Longrightarrow a^{-1}\in H.
\]
例如在整数加群中，自然数集合对加法闭合却不对取逆闭合；非零整数集合对取加法逆元闭合，却不对加法闭合。

有限群不能与其真子群同构；无限群则可以。例如整数加群与偶数加群 $2\mathbb Z$ 同构。

非交换群的真子群也可能是交换群。例如 $S_3$ 是非交换群，而
\[
H=\{(1),(12)\}=((12))
\]
是二元循环群。

一个群的两个不同子集也可能生成同一个子群。例如在 $S_3$ 中
\[
H_1=\{(123)\},\qquad H_2=\{(132)\}
\]
生成的子群都是
\[
\{(1),(123),(132)\}.
\]

\subsection*{2. 不变子群}

\textbf{子群的陪集。} 若 $H\le G$，规定
\[
a\sim b\Longleftrightarrow ab^{-1}\in H.
\]
这是等价关系，其等价类是 $H$ 的右陪集，记为 $Ha$。类似地，左陪集记作 $aH$。右（或左）陪集的个数称为 $H$ 在 $G$ 中的\textbf{指数}。

\textbf{不变子群。} 子群 $N\le G$ 若对每个 $a\in G$ 都有
\[
Na=aN,
\]
则称 $N$ 为\textbf{不变（正规）子群}。

一般子群的左、右陪集未必一致。例如 $S_3$ 中
\[
H=\{(1),(12)\}
\]
的左陪集可取
\[
(1)H=\{(1),(12)\},\quad
(13)H=\{(13),(132)\},\quad
(23)H=\{(23),(123)\},
\]
而右陪集可取
\[
H(1)=\{(1),(12)\},\quad
H(13)=\{(13),(123)\},\quad
H(23)=\{(23),(132)\}.
\]
故 $H$ 不是正规子群。

\textbf{Lagrange 定理。} 有限群 $G$ 的子群 $H$ 的阶整除 $G$ 的阶。但反之不真：若 $m\mid |G|$，$G$ 未必有阶为 $m$ 的子群。例如 $A_4$ 的阶为 $12$，但没有阶为 $6$ 的子群。

阶为素数的群一定是循环群，反之不真。例如模 $6$ 的剩余类加群是六阶循环群。

陪集相等也不推出代表元相等。例如在 $S_3$ 中取
\[
H=\{(1),(123),(132)\},
\]
有
\[
H(23)=H(13),
\]
但 $(23)\ne(13)$。

\subsection*{3. 子集的积}

若 $S_1,S_2,\ldots,S_m$ 是群 $G$ 的子集，定义
\[
S_1S_2\cdots S_m
=\{s_1s_2\cdots s_m\mid s_i\in S_i\}.
\]
若 $H$ 是子群，则 $HH=H$；反之不真。例如 $\mathbb Q^*$ 在乘法下为群，取 $H$ 为全体非零奇整数，则 $HH=H$，但 $H$ 不是群，因为例如 $3^{-1}=1/3\notin H$。

两个子群的积也未必是子群。例如在 $S_3$ 中取
\[
H_1=((12)),\qquad H_2=((23)),
\]
则
\[
H_1H_2=\{(1),(23),(12),(132)\}
\]
不构成子群。

正规性的传递性也不成立。取 $G=S_4$，
\[
N=\{(1),(12)(34),(13)(24),(14)(23)\},
\qquad
N_1=\{(1),(14)(23)\}.
\]
$N\triangleleft S_4$，且因为 $N$ 交换，有 $N_1\triangleleft N$；但
\[
(34)^{-1}(14)(23)(34)=(13)(24)\notin N_1,
\]
所以 $N_1$ 不是 $S_4$ 的正规子群。

交换群的每个子群当然都是正规子群；反之不真。存在非交换群而所有子群均正规，这类群称为\textbf{Hamilton 群}。例如在 $S_8$ 中令
\[
\pi_1=(1234)(5678),
\qquad
\pi_2=(1537)(2846),
\]
并令 $K=(\pi_1,\pi_2)$。直接计算得
\[
\pi_1^4=\pi_2^4=e,
\qquad
\pi_1^2=\pi_2^2,
\qquad
\pi_1\pi_2\pi_1=\pi_2,
\]
且
\[
K=\{e,\pi_1,\pi_2,\pi_1\pi_2,\pi_2\pi_1,\pi_1^2,\pi_1^3,\pi_2^3\}.
\]
其中
\begin{align*}
\pi_1\pi_2&=(1836)(2745),\\
\pi_2\pi_1&=(1638)(2547),\\
\pi_1^2=\pi_2^2&=(13)(24)(57)(68),\\
\pi_1^3&=(1432)(5876),\\
\pi_2^3&=(1735)(2648).
\end{align*}
因此 $K$ 是八阶非交换群。它只有一个二阶子群
\[
N_1=(\pi_1^2),
\]
以及三个四阶子群
\[
N_2=(\pi_1),\qquad N_3=(\pi_2),\qquad N_4=(\pi_1\pi_2),
\]
这些子群都正规，所以 $K$ 是 Hamilton 群。

右陪集集合与左陪集集合之间总有一一映射
\[
Ha\longmapsto a^{-1}H.
\]
但看似更直接的对应 $Ha\mapsto aH$ 未必是良定义的。仍在 $S_3$ 中取 $H=\{(1),(12)\}$，有
\[
H(13)=H(123),
\]
而
\[
(13)H\ne(123)H.
\]

最后，一个交换群的每个子群都是正规子群，但不能把“交换群”换成“交换子群”。例如 $S_3$ 中 $H=\{(1),(12)\}$ 本身是交换群，却不是 $S_3$ 的正规子群，因为
\[
(132)(12)(123)=(23)\notin H.
\]

\section{商群}

\textbf{定义。} 群 $G$ 的正规子群 $N$ 的陪集所组成的群叫做一个\textbf{商群}，记为 $G/N$。

交换群的商群仍是交换群，但反之不真：非交换群的某些商群也可以是交换群。为说明这一点，引入换位子。

\textbf{换位子。} 群 $G$ 中形如
\[
a^{-1}b^{-1}ab
\]
的元素叫做换位子。令 $C$ 为所有有限个换位子乘积组成的集合。单位元本身也是换位子，所以 $C$ 非空。换位子的有限乘积对乘法与取逆封闭，因此 $C$ 是子群；并且对任意 $c\in C,g\in G$，可由
\[
gcg^{-1}c^{-1}\in C
\]
推出 $gcg^{-1}\in C$，故 $C\triangleleft G$。

于是商群 $G/C$ 是交换群。事实上若 $x,y\in G$，则
\[
x^{-1}y^{-1}xy\in C,
\]
从而
\[
xy\in yxC,
\qquad
xyC=yxC,
\]
即
\[
(xC)(yC)=(yC)(xC).
\]

再说明集合的像与逆像。设 $\phi:A\to\bar A$ 是满射，$S\subseteq A$。$S$ 的像是所有 $\phi(s)$ 所成的集合；若 $\bar S\subseteq\bar A$，则 $\bar S$ 的逆像是所有映入 $\bar S$ 的元素。若 $S$ 是某个 $\bar S$ 的逆像，则 $\bar S$ 必是 $S$ 的像；反之不真。

例如
\[
A=\{1,2,3,4,5,6\},\qquad \bar A=\{1,2\},
\]
定义
\[
1\mapsto1,\quad2\mapsto2,\quad3\mapsto1,
\quad4\mapsto2,\quad5\mapsto1,\quad6\mapsto2.
\]
令 $S=\{1,3\}$，则 $S$ 的像为 $\{1\}$，但 $\{1\}$ 的逆像为 $\{1,3,5\}$，并不等于 $S$。

最后讨论 $G$ 与商群 $G/N$ 的同态。自然同态
\[
\phi:G\to G/N,
\qquad
x\longmapsto Nx
\]
的核正是 $N$。但若另取一个同态满射 $G\to G/N$，其核未必仍是 $N$。

例如取 Klein 四元群
\[
G=\{e,a,b,c\},
\qquad
\begin{array}{c|cccc}
 &e&a&b&c\\\hline
 e&e&a&b&c\\
 a&a&e&c&b\\
 b&b&c&e&a\\
 c&c&b&a&e
\end{array}
\]
以及正规子群
\[
N=\{e,a\}.
\]
自然同态 $x\mapsto Nx$ 的核为 $N$。另一方面定义
\[
\psi(e)=Ne,
\qquad
\psi(b)=Ne,
\qquad
\psi(c)=Nb,
\qquad
\psi(a)=Nb.
\]
容易验证 $\psi$ 也是 $G$ 到 $G/N$ 的满射同态，但
\[
\ker\psi=N_1=\{e,b\}\ne N.
\]
根据同态基本定理
\[
G/N_1\cong G/N.
\]
这同时说明：同一个群的两个商群可以同构，而相应的正规子群未必相同。
